% 分野別類題: 波動方程式（ダランベールの解）
% u_tt = c^2 u_xx（無限区間）を特性変数 xi = x - ct, eta = x + ct で解く基本形 3 問。
% コンパイル: TEXINPUTS="../../../00_common:$TEXINPUTS" latexmk -xelatex problems.tex
\documentclass[a4paper,11pt]{article}
\usepackage{preamble}

\begin{document}

\kamokumei{類題演習 --- 波動方程式（ダランベールの解）}

\shijibun{次の波動方程式（無限区間）に関する問いに答えよ。（1次元波動方程式 $u_{tt} = c^{2} u_{xx}$ は特性変数 $\xi = x - ct$, $\eta = x + ct$ への変換により一般解 $u = f(x-ct) + g(x+ct)$ をもつ。）}

\shomon{問1.}{波動方程式
\[
\frac{\partial^{2} u}{\partial t^{2}} = c^{2}\frac{\partial^{2} u}{\partial x^{2}} \qquad (-\infty < x < \infty,\ t > 0)
\]
に対し、特性変数 $\xi = x - ct$, $\eta = x + ct$ を用いて変数変換を行い、一般解が
\[
u(x,t) = f(x-ct) + g(x+ct)
\]
の形で表されることを導け（$f$, $g$ は任意の $C^{2}$ 級関数）。}

\shomon{問2.}{波動方程式 $u_{tt} = c^{2} u_{xx}$ が初期条件
\[
u(x,0) = \varphi(x), \qquad u_{t}(x,0) = 0
\]
（初期変位のみ、初期速度なし）を満たすとき、$u(x,t)$ をダランベールの公式を用いて $\varphi$ で表せ。特に $\varphi(x) = e^{-x^{2}}$、$c = 2$ のときの $u(x,t)$ を具体的に求めよ。}

\shomon{問3.}{波動方程式 $u_{tt} = c^{2} u_{xx}$ が初期条件
\[
u(x,0) = 0, \qquad u_{t}(x,0) = \psi(x)
\]
（初期速度のみ、初期変位なし）を満たすとき、$u(x,t)$ をダランベールの公式を用いて $\psi$ で表せ。特に $\psi(x) = \cos x$、$c = 3$ のときの $u(x,t)$ を具体的に求めよ。}

\end{document}
